\documentclass[11pt,letterpaper]{article}
\usepackage[margin=1in]{geometry}
\usepackage[T1]{fontenc}
\usepackage{lmodern}
\usepackage{microtype}
\usepackage{amsmath,amssymb,amsthm}
\usepackage{booktabs}
\usepackage{enumitem}
\usepackage{setspace}
\usepackage[numbers,sort&compress]{natbib}
\usepackage[hidelinks]{hyperref}
\usepackage[nameinlink,noabbrev]{cleveref}
\setstretch{1.08}
\setlength{\parindent}{1.5em}
\setlength{\parskip}{0.25em}
\newtheorem{proposition}{Proposition}
\title{\textbf{A Hybrid Lanchester--Markov Framework for Probabilistic Conflict-State Dynamics}}
\author{Krish Parikh (Lapenzer)\\Independent Researcher, Ontario, Canada\\\texttt{lapenzerparikh@gmail.com}}
\date{Preprint -- not peer reviewed\\Copyright \textcopyright\ 2026 Krish Parikh. All rights reserved.}

\begin{document}
\maketitle

\begin{abstract}
This paper develops a theoretical hybrid model coupling Lanchester-type continuous attrition dynamics with a networked discrete-state Markov model of geographic conflict evolution. Classical Lanchester equations model the continuous evolution of two opposing aggregate force quantities but do not directly represent uncertainty over strategic geographic states. Conversely, a Markov conflict model can represent probabilistic changes among states such as stable, pressured, contested, and changed control but does not inherently explain how local force attrition modifies those transition probabilities. The proposed framework joins the two representations through an augmented state process. During each discrete forecast interval, state-dependent Lanchester equations evolve continuous effective-force variables. Derived quantities---including force ratios, attrition fractions, and remaining effective strength---then enter the transition kernel of the city-level Markov model. The newly sampled geographic state determines the parameter regime governing the next continuous interval, producing a stochastic switching dynamical system. Parameter uncertainty can be propagated using Bayesian or Monte Carlo sampling rather than fixed deterministic attrition coefficients. The paper establishes the Markov property of the augmented process under explicit assumptions and introduces stopping conditions that prevent physically meaningless negative force values. The framework deliberately treats Lanchester dynamics as one information source rather than a deterministic law of war. No empirical validation, historical fit, or operational predictive accuracy is claimed.
\end{abstract}

\noindent\textbf{Keywords:} Lanchester equations; Markov chains; stochastic hybrid systems; Monte Carlo simulation; attrition; conflict forecasting; Bayesian inference

\section{Introduction}
A city-state Markov model can answer questions of the form
\begin{equation}
P(S_{i,t+1}=s'\mid\mathbf{S}_t),
\end{equation}
but it does not, by itself, describe how changes in opposing force strength affect those transition probabilities. Lanchester-type equations provide one possible continuous representation of aggregate attrition.

The central question of this paper is therefore whether continuous Lanchester attrition dynamics can be embedded inside a probabilistic Markov conflict model without treating either representation as deterministically correct. The proposed architecture contains both continuous force dynamics and discrete probabilistic geographic states:
\begin{equation}
\text{Lanchester dynamics}\rightarrow\text{Markov transition features}\rightarrow\text{new conflict state}\rightarrow\text{new Lanchester regime}.
\end{equation}
Neither representation replaces the other.

\section{Related Mathematical Work}
Lanchester models have a long history as aggregate mathematical models of attrition. Deterministic formulations are widely discussed in operations-research and mathematical-combat literature. Stochastic extensions have also been developed. Kearney and Martin examined a stochastic version of Lanchester combat dynamics \citep{kearney2019}. Pettit, Wiper, and Young developed Bayesian inference for stochastic Lanchester laws \citep{pettit2003}, and Kim et al. studied approximation methods for stochastic Lanchester models \citep{kim2017}.

Markov representations have separately been used for conflict escalation and de-escalation forecasting \citep{randahl2022}. The present work proposes a specific coupling:
\begin{equation}
\boxed{\text{Lanchester continuous state}\leftrightarrow\text{network Markov discrete state}}
\end{equation}
with Monte Carlo propagation of both transition and attrition uncertainty.

\section{Discrete Geographic State}
Let $G=(V,E,W)$ be a city network. For each city $i$,
\begin{equation}
S_{i,t}\in\mathcal{S}=\{0,1,2,3\},
\end{equation}
where 0 denotes stable, 1 pressured, 2 contested, and 3 control changed. Let
\begin{equation}
\mathbf{S}_t=(S_{1,t},\ldots,S_{n,t}).
\end{equation}

\section{Continuous Force State}
For each relevant city or geographic sector $i$, define two non-negative effective-force quantities
\begin{equation}
F_i(\tau)\geq0,
\qquad
G_i(\tau)\geq0.
\end{equation}
These need not represent literal troop counts; they may represent normalized effective strength, provided the definition remains consistent.

During a local forecast interval, the classical aimed-fire Lanchester structure is
\begin{equation}
\frac{dF_i}{d\tau}=-\beta_iG_i,
\qquad
\frac{dG_i}{d\tau}=-\alpha_iF_i,
\label{eq:lanchester}
\end{equation}
where $\alpha_i,\beta_i\geq0$ are effectiveness parameters.

\section{Classical Invariant}
For constant $\alpha_i$ and $\beta_i$, consider
\begin{equation}
Q_i=\alpha_iF_i^2-\beta_iG_i^2.
\end{equation}
Differentiating and substituting \cref{eq:lanchester} gives
\begin{align}
\frac{dQ_i}{d\tau}
&=2\alpha_iF_i\frac{dF_i}{d\tau}-2\beta_iG_i\frac{dG_i}{d\tau}\\
&=2\alpha_iF_i(-\beta_iG_i)-2\beta_iG_i(-\alpha_iF_i)=0.
\end{align}
Therefore
\begin{equation}
\alpha_iF_i^2-\beta_iG_i^2=\text{constant}
\label{eq:invariant}
\end{equation}
within a constant-parameter interval. This is a property of the idealized mathematical model, not a universal empirical law of actual warfare.

\section{Stopping Rule}
Unconstrained continuous solutions can become physically meaningless if mathematically extended beyond elimination of one force. Define
\begin{equation}
\tau_i^{\star}=\inf\{\tau\geq0:F_i(\tau)=0\ \text{or}\ G_i(\tau)=0\}.
\end{equation}
The continuous integration interval is
\begin{equation}
[0,\min(\Delta,\tau_i^{\star})].
\end{equation}
After the stopping time, non-negativity is preserved through an absorbing boundary at zero unless a separate reinforcement model is introduced.

\section{State-Dependent Lanchester Parameters}
Attrition dynamics need not use identical parameters in every geographic state. Define
\begin{equation}
\alpha_i=\alpha_i(S_{i,t}),
\qquad
\beta_i=\beta_i(S_{i,t}).
\end{equation}
Thus a transition from one discrete geographic state to another may change the parameter regime used during the next interval. More generally,
\begin{equation}
\alpha_{i,t}=f_{\alpha}(S_{i,t},\mathbf{x}_{i,t},\theta_{\alpha}),
\qquad
\beta_{i,t}=f_{\beta}(S_{i,t},\mathbf{x}_{i,t},\theta_{\beta}).
\end{equation}
The paper does not assert what those coefficients should be; they require empirical estimation.

\section{Lanchester-Derived Features}
After integrating over a forecast interval $\Delta$, define
\begin{equation}
F_{i,t}^{+}=F_i(\Delta),
\qquad
G_{i,t}^{+}=G_i(\Delta).
\end{equation}
A log effective-force ratio is
\begin{equation}
R_{i,t}=\log\left(\frac{F_{i,t}^{+}+\epsilon}{G_{i,t}^{+}+\epsilon}\right).
\end{equation}
Attrition fractions are
\begin{equation}
L_{F,i,t}=\frac{F_{i,t}-F_{i,t}^{+}}{F_{i,t}+\epsilon},
\qquad
L_{G,i,t}=\frac{G_{i,t}-G_{i,t}^{+}}{G_{i,t}+\epsilon}.
\end{equation}
Define the continuous feature vector
\begin{equation}
\mathbf{L}_{i,t}=
\begin{bmatrix}
R_{i,t} & L_{F,i,t} & L_{G,i,t} & F_{i,t}^{+} & G_{i,t}^{+}
\end{bmatrix}^{\top}.
\end{equation}
These variables connect the continuous and discrete parts of the model.

\section{Markov Transition Kernel with Attrition}
For candidate next state $s'$, define
\begin{equation}
\eta_{i,t}(s')
=
a_{S_{i,t},s'}
+\boldsymbol{\beta}_{s'}^{\top}\mathbf{x}_{i,t}
+\boldsymbol{\gamma}_{s'}^{\top}\mathbf{L}_{i,t}
+\sum_{j\in N(i)}\omega_{ji,s'}\phi(S_{j,t}).
\end{equation}
Then
\begin{equation}
P(S_{i,t+1}=s'\mid\mathbf{S}_t,\mathbf{F}_t,\mathbf{G}_t,\mathbf{X}_t)
=
\frac{\exp(\eta_{i,t}(s'))}{\sum_{r\in\mathcal{S}}\exp(\eta_{i,t}(r))}.
\label{eq:hybrid-transition}
\end{equation}
Consequently, Lanchester outputs influence probabilities rather than deterministically deciding the next state. Even a large force advantage does not mathematically force a state transition with probability one.

\section{Augmented State}
Define
\begin{equation}
Z_t=(\mathbf{S}_t,\mathbf{F}_t,\mathbf{G}_t,\mathbf{X}_t).
\label{eq:aug-state}
\end{equation}
The transition from $Z_t$ to $Z_{t+1}$ consists of two stages. First, the continuous state evolves according to the selected Lanchester regime. Second, the discrete state is sampled from \cref{eq:hybrid-transition}. The newly sampled state then determines the parameter regime for the next continuous interval.

\begin{proposition}[Markov property of the hybrid process]
Suppose: (i) continuous evolution over period $t$ depends only on $Z_t$ and a current-period innovation; (ii) innovations are conditionally independent of earlier states given $Z_t$; (iii) the discrete transition kernel depends only on the post-flow continuous state and current discrete state; and (iv) exogenous covariates are included in $Z_t$ or follow a Markov process included in the augmented state. Then $\{Z_t\}$ is a Markov process.
\end{proposition}
\begin{proof}
The continuous state at the end of period $t$ is determined conditionally by the present augmented state and current innovation. The discrete transition depends only on that continuous result and the present state. Therefore
\begin{equation}
P(Z_{t+1}\mid Z_0,\ldots,Z_t)=P(Z_{t+1}\mid Z_t),
\end{equation}
which is the Markov property.
\end{proof}

The discrete geographic state alone need not be Markov if the continuous force variables are omitted.

\section{Why the Augmented State Matters}
Suppose two scenarios have identical geographic states but different effective-force conditions:
\begin{equation}
\mathbf{S}_t^{(1)}=\mathbf{S}_t^{(2)},
\qquad
\mathbf{F}_t^{(1)}\neq\mathbf{F}_t^{(2)}.
\end{equation}
Then their next-state distributions may differ. The apparent violation of the Markov assumption disappears when force state is included in the augmented state $Z_t$. Lanchester variables can therefore be interpreted as part of the state information required to make the process approximately Markovian.

\section{Stochastic Attrition Coefficients}
Fixed coefficients can give a misleading appearance of certainty. Instead let
\begin{equation}
\alpha_{i,t}\sim p_{\alpha}(\alpha\mid D_t),
\qquad
\beta_{i,t}\sim p_{\beta}(\beta\mid D_t).
\end{equation}
A Monte Carlo realization therefore samples both attrition parameters and discrete state transitions. For trial $m$, sample
\begin{equation}
\alpha_{i,t}^{(m)}\sim p_{\alpha},
\qquad
\beta_{i,t}^{(m)}\sim p_{\beta},
\end{equation}
then integrate the Lanchester system using those sampled parameters and feed the resulting $\mathbf{L}_{i,t}^{(m)}$ into the Markov transition kernel.

\section{Bayesian Formulation}
Let $\Theta_L$ contain Lanchester parameters and $\Theta_M$ contain Markov parameters. Given observations $D$,
\begin{equation}
p(\Theta_L,\Theta_M\mid D)
\propto
p(D\mid\Theta_L,\Theta_M)p(\Theta_L,\Theta_M).
\end{equation}
Posterior predictive simulation becomes
\begin{equation}
p(Z_{t+1:T}\mid Z_t,D)
=
\int p(Z_{t+1:T}\mid Z_t,\Theta)p(\Theta\mid D)\,d\Theta.
\end{equation}
Monte Carlo approximates this integral by sampling parameter vectors from the posterior and then generating future trajectories.

\section{Network Interaction}
Local force conditions should not erase geographic cascade structure. For city $i$, define
\begin{equation}
C_{i,t}=\sum_{j\in N(i)}w_{ji}\phi(S_{j,t}).
\end{equation}
A generic binary transition equation is
\begin{equation}
\operatorname{logit}P(S_{i,t+1}=3)
=
\theta_0+\theta_RR_{i,t}+\theta_FL_{F,i,t}+\theta_GL_{G,i,t}+\theta_CC_{i,t}+\boldsymbol{\theta}_X^{\top}\mathbf{x}_{i,t}.
\label{eq:binary-hybrid}
\end{equation}
This is an estimable statistical specification, not a statement that its coefficients are presently known.

\section{Feedback}
The model contains a feedback loop:
\begin{equation}
S_{i,t}
\rightarrow
(\alpha_{i,t},\beta_{i,t})
\rightarrow
(F_{i,t}^{+},G_{i,t}^{+})
\rightarrow
P(S_{i,t+1})
\rightarrow
S_{i,t+1}.
\end{equation}
Consequently, a discrete state change can alter future continuous dynamics, and continuous dynamics can alter future state probabilities.

\section{Monte Carlo Algorithm}
\begin{enumerate}[leftmargin=2em]
\item Initialize geographic states, effective-force variables, network, contextual covariates, parameter distributions, horizon $T$, and simulation count $M$.
\item For each simulation and each period, sample uncertain Lanchester parameters.
\item Integrate the continuous attrition equations until time $\Delta$ or an elimination boundary.
\item Compute Lanchester-derived and network-state features.
\item Evaluate Markov transition probabilities and sample each city's next state.
\item Update state-dependent parameter regimes and advance to the next period.
\item Store each complete trajectory.
\item Estimate state probabilities, cascade probabilities, distributions of remaining effective strength, and simulation intervals.
\item Report sensitivity to parameters and assumptions.
\end{enumerate}

\section{Illustrative Two-City Structure}
Consider two cities $A$ and $B$. Suppose Lanchester integration in A produces force feature vector $\mathbf{L}_{A,t}$. Then a local transition probability may take the form
\begin{equation}
P(S_{A,t+1}=3)=\sigma(\theta_A^{\top}\mathbf{L}_{A,t}+\boldsymbol{\beta}_A^{\top}\mathbf{x}_{A,t}).
\end{equation}
If the resulting probability is 0.50, and
\begin{equation}
P(S_{B,t+2}=3\mid S_{A,t+1}=3)=q_1,
\qquad
P(S_{B,t+2}=3\mid S_{A,t+1}\neq3)=q_0,
\end{equation}
then
\begin{equation}
P(S_{B,t+2}=3)=0.5q_1+0.5q_0.
\end{equation}
Thus the network propagation equation naturally nests inside the hybrid model. The difference is that A's transition probability may now depend partly on the continuous attrition state.

\section{Model Hierarchy}
The framework can be interpreted as three nested layers. Layer 1 maps effective force to an attrition trajectory. Layer 2 maps attrition trajectory and context to a local city-state transition distribution. Layer 3 maps city-state distributions into regional cascade distributions. This prevents any single layer from being interpreted as a complete theory of conflict.

\section{Evaluation}
No validation is presented. A future empirical study would need to determine whether Lanchester-type equations adequately approximate the selected level of aggregation, whether their parameters can be estimated reliably, whether Lanchester-derived variables improve transition forecasts beyond a Markov baseline, whether the network component adds predictive information, and whether forecast probabilities are calibrated.

A key model hierarchy is
\begin{equation}
M_0=\text{Markov only},\qquad
M_1=\text{Markov + network},\qquad
M_2=\text{Markov + network + Lanchester}.
\end{equation}
The more complicated model should not be preferred unless out-of-sample evidence justifies the additional complexity.

\section{Limitations of Lanchester Dynamics}
Lanchester equations make strong simplifications. They aggregate heterogeneous forces and do not automatically represent terrain, logistics, command quality, morale, intelligence, fortification, civilian behaviour, technological asymmetry, political constraints, reinforcement, or tactical geometry. Continuous attrition can also be inappropriate when events are sparse or highly discrete. Effectiveness coefficients may not remain constant.

For these reasons, the proposed framework does not use Lanchester equations as deterministic predictors of geographic outcomes. Instead, Lanchester output is treated as one probabilistic input to a broader transition model.

\section{Model Risk}
The hybrid model contains at least three forms of uncertainty: state uncertainty, parameter uncertainty, and structural uncertainty. Monte Carlo methods can represent the first two more readily than the third. Structural uncertainty therefore requires explicit comparison with alternative models.

\section{Contribution}
The proposed contribution is the architecture
\begin{equation}
\begin{aligned}
&\text{networked Markov state}+\text{Lanchester continuous state}+\text{state-dependent switching}\\
&\qquad+\text{Bayesian uncertainty}+\text{Monte Carlo propagation}
\end{aligned}
\end{equation}
rather than any claim to have invented the individual components. The key theoretical move is to enlarge the Markov state from $\mathbf{S}_t$ to the augmented state in \cref{eq:aug-state}, thereby allowing continuous attrition variables to preserve information relevant to future discrete transitions.

\section{Conclusion}
This paper proposed a stochastic hybrid system coupling Lanchester-type attrition dynamics with a networked Markov model of geographic conflict states. The continuous layer produces distributions over changing effective-force conditions. The discrete layer translates those conditions, together with contextual and network variables, into probabilities over future geographic states. The newly sampled state then modifies the next interval's attrition regime.

The framework is presently theoretical. Its value cannot be established from mathematical coherence alone. Empirical research would be required to determine whether Lanchester-derived variables add forecast information beyond simpler Markov and network models.

\section*{Statement on AI Assistance}
OpenAI ChatGPT was used as an editorial, mathematical-formalization, and literature-organization assistant during preparation of this manuscript. The author retains responsibility for the research concepts, interpretation, verification, and final manuscript.

\begin{thebibliography}{99}
\bibitem{kearney2019} Kearney, M. J., \& Martin, R. J. (2019). On a stochastic version of Lanchester's model of combat. \emph{arXiv preprint arXiv:1905.03122}.
\bibitem{kim2017} Kim, D., Moon, H., Park, D., \& Shin, H. (2017). An efficient approximate solution for stochastic Lanchester models. \emph{Journal of the Operational Research Society}, 68(11), 1470--1481. doi:10.1057/s41274-016-0163-6.
\bibitem{pettit2003} Pettit, L. I., Wiper, M. P., \& Young, K. D. S. (2003). Bayesian inference for some Lanchester combat laws. \emph{European Journal of Operational Research}, 148(1), 152--165. doi:10.1016/S0377-2217(02)00462-9.
\bibitem{randahl2022} Randahl, D., \& Vegelius, J. (2022). Predicting escalating and de-escalating violence in Africa using Markov models. \emph{International Interactions}, 48(4), 597--613. doi:10.1080/03050629.2022.2049772.
\end{thebibliography}
\end{document}
